\chapter*{Preface}
\addcontentsline{toc}{chapter}{Preface}

A semiconductor course can become a list of formulas remarkably quickly.  A
band gap appears, two effective densities of states are quoted, a Fermi level
is placed somewhere in the gap, and a concentration falls out of an
exponential.  The calculation may be correct, yet the connection between the
crystal, the model, and the reported number remains hidden.

These notes take the opposite route.  Each important quantity is attached to a
model and each model is attached to evidence.  The derivation identifies the
origin of the band energies, the range replaced by a parabola, the information
preserved by an effective mass, the multiplicity of the available states, the
partition function governing occupation, and the charge constraint determining
the chemical potential.  Computation provides an inspectable representation of
the physical and mathematical argument.

The recurring modeling chain is
\begin{equation*}
  H(\bm k)
  \longrightarrow E_n(\bm k)
  \longrightarrow \text{state multiplicity}
  \longrightarrow g(E)
  \longrightarrow Z,\Xi
  \longrightarrow f(E,\mu,T)
  \longrightarrow n,p
  \longrightarrow Q(\mu)=0.
\end{equation*}
Every arrow carries assumptions, units, conventions, and a range of validity.
Semiconductor physics requires selecting reductions that remain accurate in the
regime under study.

Each treatment states the physical model and its domain, gives an inspectable
analytic or numerical representation, and supplies validation evidence tied to
the defining physical relations.

The initial chapters connect the particle in a box to Bloch states and finite
HOMO--LUMO spectra to tight-binding bands.  They then establish spectral
multiplicity, density of states, partition functions, occupation, chemical
potential, and equilibrium.  The intrinsic and extrinsic semiconductor chapters
then apply
that foundation before later treatments of transport, recombination,
interfaces, and junctions.  Each treatment uses the simplest explicit model
that preserves the physics required by the problem.

The principal textbook conventions follow standard treatments of solid-state
and semiconductor physics, including Refs.~\cite{kittel2005,neamen2012,yu2010}
and the device-oriented perspectives of Refs.~\cite{szeLee2012,szeNg2007}.  Specific material
claims require material-specific primary or curated sources with declared
provenance.
